\chapter*{PHỤ LỤC}
\addcontentsline{toc}{chapter}{Phụ lục}

\section*{Phụ lục 1: Danh mục các tính năng đã chuẩn hóa}

\textit{Hệ thống hiện đã hoàn thiện việc chuẩn hóa giao diện theo shadcn/ui, bao gồm Bảng điều khiển (Dashboard), Cộng đồng (Social Hub), Trung tâm tin nhắn (Chat) và các công cụ rèn luyện kỷ luật (Grid, Journal).}

%% ────────────────────────────────────────────
\section*{Phụ lục 2: Hướng dẫn nhanh dành cho Mentor (Quick Start Guide)}

\begin{enumerate}
  \item \textbf{Đăng nhập}: Truy cập địa chỉ nội bộ, sử dụng Email cá nhân. Nếu quên mật khẩu, hãy dùng tính năng ``Forgot Password''.
  \item \textbf{Giao diện chính}:
    \begin{itemize}
        \item \textbf{Cột trái}: Menu điều hướng (Dashboard, Mentees, Goals, Meetings, Wiki).
        \item \textbf{Trung tâm}: Các thẻ chỉ số (Stats) và Lịch họp sắp tới.
    \end{itemize}
  \item \textbf{Xem chi tiết Mentee}: Click vào danh sách Mentee. 
    \begin{itemize}
        \item Quan sát \textbf{Heatmap}: Tìm kiếm những ô màu xám (ngày không rèn luyện).
        \item Quan sát \textbf{Goal Progress}: Kiểm tra xem thanh \% có di chuyển trong 7 ngày qua không.
    \end{itemize}
  \item \textbf{Ghi chép Minutes}: Trong buổi họp, click vào ``New Meeting'' để ghi lại các ý chính. Điều này giúp hệ thống tự động nhắc nhở Mentee làm Reflection sau đó.
\end{enumerate}

%% ────────────────────────────────────────────
\section*{Phụ lục 3: Checklist Phản tư (Reflection) dành cho Mentee}

Để đạt hiệu quả cao nhất từ WebApp, Mentee nên trả lời 5 câu hỏi sau trong bài Reflection hằng tuần:
\begin{itemize}
    \item \textbf{Q1 - Chiến thắng}: Điều gì em đã làm tốt nhất tuần vừa qua?
    \item \textbf{Q2 - Thử thách}: Khó khăn lớn nhất em gặp phải là gì? Tại sao?
    \item \textbf{Q3 - Bài học}: Mentor đã chỉ dạy điều gì khiến em thay đổi tư duy?
    \item \textbf{Q4 - Hành động}: Tuần tới em sẽ thay đổi cụ thể điều gì trong bộ Task của mình?
    \item \textbf{Q5 - Yêu cầu}: Em cần Mentor hỗ trợ thêm nguồn lực/tri thức gì?
\end{itemize}

%% ────────────────────────────────────────────
\section*{Phụ lục 4: Giải quyết vấn đề dành cho Mentor cao tuổi (FAQ)}

\begin{table}[H]
    \centering
    \begin{tabularx}{\textwidth}{|L|L|}
    \hline
    \rowcolor{pblue} \textbf{\color{white} Câu hỏi} & \textbf{\color{white} Giải pháp} \\ \hline
    Tôi thấy chữ hơi nhỏ & Nhấn tổ mã phím \texttt{Ctrl +} (Windows) hoặc \texttt{Cmd +} (Mac) để phóng to trình duyệt. \\ \hline
    Tôi không biết Mentee đã viết bài chưa & Vào menu ``Reflections'', hệ thống sẽ liệt kê các bài mới nhất của tất cả Mentee lên đầu. \\ \hline
    Tôi muốn nhập dữ liệu thay Mentee & Hệ thống khuyến khích Mentee tự nhập để rèn tính kỷ luật, nhưng Mentor vẫn có quyền Edit các Task nếu cần hướng dẫ cụ thể. \\ \hline
    WebApp có dùng được trên điện thoại không? & Có, WebApp được thiết kế Responsive để hiển thị tốt trên iPhone/Android thông qua trình duyệt web. \\ \hline
    \end{tabularx}
    \caption{Các câu hỏi thường gặp của Mentor khi mới tiếp cận WebApp}
    \captionsource{Tulie Tech Research, 2026}
\end{table}

%% ────────────────────────────────────────────
\section*{Phụ lục 5: Chân dung người dùng (User Personas)}

\begin{itemize}
    \item \textbf{Mentor A (Bận rộn)}: Chủ doanh nghiệp, chỉ có 45 phút/tuần cho Mentee. WebApp giúp anh lướt Report trong 2 phút trước họp $\to$ Đi thẳng vào coaching.
    \item \textbf{Mentee B (Cần kỷ luật)}: Sinh viên năm cuối, năng lượng thất thường. WebApp và Streak Counter giúp bạn tự tạo áp lực tích cực để không bỏ cuộc.
    \item \textbf{Ban Điều Hành}: Cần báo cáo cuối kỳ cho Viện ISME. WebApp hỗ trợ xuất báo cáo dữ liệu sạch (clean data) trong 1 click.
\end{itemize}
